本研究はPubMedという医学生物学分野に限られた範囲のデータにおける分析であり、他の分野あるいは複数の分野を対象とした場合に同様の結果と考察を得られるかは分からない。

一つの検証として、前述の最多被引用論文2本につき、本研究による被引用数と自然科学領域全体を対象とした"The top 100 papers"によるそれを比較（本研究の1946-2015累積世代とWeb of Scienceの2014年データ）すると、
"Protein measurement with the Folin phenol regent"では、前者で49164件および後者で305148件、
"Cleavage of structural protein during the assembly of the head of bacteriophage T4"では、前者で49916件および後者で213005件
であり、およそ4倍から6倍の差があり、同様の解釈はできない。
とはいえ、本研究の手法でも学際領域における生命科学・生物科学分野での上位の論文を抽出できており、順位尺度としては概ね対照的な関係となっている可能性が高い。

雑誌レベルの観点では、"The top 100 papers"においても必ずしも高インパクトファクター誌が現れるとは限らないことが示されており、この点においては学際的領域においても医学生物学分野と同様の性質があると考えられる。

一方、研究テーマの分析に関しては、現実的に利用可能な学祭的なデータベース（Web of Science・OpenAlex等）のすべての分野において、本研究と同等の粒度および付与規則のディスクリプターを得ることは難しく、この再現性を維持したまま対象分野を広げることは難しい。

以上を勘案した今後の課題としては、(1)異なる分野においても同様の性質が見られるか、(2)学際的な領域においても同様の性質が見られるか、がある。また、この際にテーマ分析においてはその比較が困難であることが予想され、単純な領域変更・拡大の観点からはスコープ外となるが、新たなテーマとして、(3)研究テーマのマッピング、を考えることができる。
